\begin{table}[!ht]
\centering
\begin{threeparttable}
\caption{Directional Accuracy by Model and Prediction Horizon}
\label{tab:eval_da}

\begin{tabular}{lrrrrrr}
\toprule
Model & DA 1d & DA 3d & DA 5d & Naive 1d & Naive 3d & Naive 5d\\
\midrule
OLS & 0.3219 & 0.4762 & 0.4990 & 0.2762 & 0.32 & 0.3029\\
Random Forest & 0.3409 & 0.5010 & 0.4724 & 0.2762 & 0.32 & 0.3029\\
XGBoost & 0.2857 & 0.5219 & 0.4419 & 0.2762 & 0.32 & 0.3029\\
SVM & 0.3200 & 0.4914 & 0.4762 & 0.2762 & 0.32 & 0.3029\\
ANN & 0.3429 & 0.4648 & 0.4762 & 0.2762 & 0.32 & 0.3029\\
\bottomrule
\end{tabular}
\begin{tablenotes}
\footnotesize
\item \textit{Note.} DA = fraction of test-set observations where predicted sign matches realised sign. Naive benchmark uses the previous day BPI log return (lag1). Values above 0.50 indicate better-than-random directional prediction.
\end{tablenotes}
\end{threeparttable}
\end{table}
